%%% BEGIN fig:coleman-norton
\begin{figure}[htbp]
\centering
\begin{tikzpicture}[fd,
  t2 cone/.style={draw=gray!60,line width=0.5pt,dash pattern=on 3pt off 2pt},
  t2 soft/.style={draw=green!45!black,line width=0.85pt,decorate,
      decoration={fd photon,amplitude=0.9mm,segment length=17mm}},
  t2 lab/.style={font=\footnotesize}]
  % axes and light cone  (x^\pm=(t\pm z)/\sqrt2)
  \draw[lcaxis] (-3.9,0) -- (4.0,0) node[below] {$z$};
  \draw[lcaxis] (0,-1.45) -- (0,3.85) node[left] {$t$};
  \draw[t2 cone] (0,0) -- (3.6,3.6) node[above right=-2pt,text=gray!50!black,font=\footnotesize] {$x^-=0$};
  \draw[t2 cone] (0,0) -- (-3.6,3.6) node[above left=-2pt,text=gray!50!black,font=\footnotesize] {$x^+=0$};
  % incoming virtual photon at rest (world line along t)
  \draw[photon] (0,-1.3) -- (0,-0.17);
  \node[t2 lab,anchor=west] at (0.1,-0.95) {$\gamma^*$（静止）};
  % jet 1: quark along x^-=0, with a collinear gluon emitted at P1
  \draw[quark=0.22] (0.11,0.11) -- (3.15,3.15);
  \coordinate (P1) at (1.95,1.95);
  \draw[gluon] (P1) -- (2.78,3.22);
  \node[vertex] at (P1) {};
  % jet 2: antiquark along x^+=0 (fermion-number arrow towards the vertex), collinear gluon at P2
  \draw[quark=0.8] (-3.15,3.15) -- (-0.11,0.11);
  \coordinate (P2) at (-1.95,1.95);
  \draw[gluon] (P2) -- (-2.78,3.22);
  \node[vertex] at (P2) {};
  % soft (zero-momentum, long-wavelength) gluon exchanged between the jets
  \coordinate (S1) at (1.3,1.3); \coordinate (S2) at (-1.25,1.25);
  \draw[t2 soft] (S2) -- (S1);
  \node[vertex] at (S1) {}; \node[vertex] at (S2) {};
  \node[t2 lab,text=green!35!black,anchor=south] at (-0.66,1.38) {软胶子};
  % hard vertex
  \filldraw[hard blob] (0,0) circle[radius=0.16];
  \node[t2 lab,anchor=north west] at (0.15,-0.05) {$H$};
  % jet labels
  \node[t2 lab,anchor=west,align=left] at (2.2,1.1) {$q$ 喷注：$\Delta x^\mu=\alpha\,p_1^\mu$\\（沿光锥 $x^-=0$）};
  \node[t2 lab,anchor=east,align=right] at (-2.2,1.1) {$\bar q$ 喷注：$\Delta x^\mu=\alpha\,p_2^\mu$\\（沿光锥 $x^+=0$）};
\end{tikzpicture}
\caption{Coleman--Norton 定理的时空图像（$e^+e^-$ 湮灭，$\gamma^*$ 静止系，$z$ 轴取为喷注轴，$x^\pm=(t\pm z)/\sqrt2$）。钳夹面上每条在壳内线 $i$
对应一段经典自由传播 $\Delta x_i^\mu=\alpha_i l_i^\mu$（$\alpha_i\ge0$，与该线的 Feynman 参数成正比）；由 Landau 方程，各顶点的时空位置彼此相容。离壳（硬）线的 $\alpha_i=0$，全部收缩为原点处的硬顶点 $H$；
有限动量的无质量共线粒子沿光锥 $x^-=0$（$q$ 喷注）或 $x^+=0$（$\bar q$ 喷注）以光速运动，并可在同一光线上任意处分裂或重新结合（图中共线胶子与夸克线的夹角是为了看清而夸大的）。
动量为零的软线（绿色长波长线）不受此约束：它在 Landau 方程中的贡献 $\alpha_i l_i^\mu$ 恒为零，与 $\alpha_i$ 无关；物理上它是波长无限长、不能局域化的量子，可以在两个喷注的任意位置被发射或吸收。
绿线只是示意：它连接的两点一般是类空间隔的。}
\label{fig:coleman-norton}
\end{figure}
%%% END fig:coleman-norton

%%% BEGIN fig:ir-mode-exponents
\begin{figure}[htbp]
\centering
\begin{tikzpicture}[fd,x={(1.9cm,0cm)},y={(-0.996cm,0.575cm)},z={(0cm,1.55cm)},
  fb ax/.style={lcaxis,gray!50!black},
  fb guide/.style={draw=gray!60,line width=0.4pt,dash pattern=on 1pt off 1.6pt},
  fb plane/.style={fill=gray!22,fill opacity=0.62,draw=gray!55!black,line width=0.6pt},
  fb pt/.style={circle,draw=white,line width=0.4pt,inner sep=0pt,minimum size=5.6pt},
  fb gl/.style={circle,draw=red!60!black,fill=white,line width=0.9pt,inner sep=0pt,minimum size=6pt},
  fb lab/.style={font=\footnotesize,inner sep=2pt},
  fb tick/.style={font=\scriptsize,text=gray!35!black,inner sep=1.5pt}]
  % ---------------- base grid (c = 0) and axes ----------------
  \foreach \t in {1,2} {
    \draw[fb guide] (\t,0,0) -- (\t,2,0);
    \draw[fb guide] (0,\t,0) -- (2,\t,0);
  }
  \draw[fb ax] (0,0,0) -- (2.75,0,0) node[right,font=\small] {$a$};
  \draw[fb ax] (0,0,0) -- (0,2.85,0) node[left,font=\small] {$b$};
  \foreach \t in {1,2} {
    \draw[gray!50!black,line width=0.5pt] (\t,0,0) -- ++(0,0,-0.07) node[below,fb tick] {$\t$};
    \draw[gray!50!black,line width=0.5pt] (0,\t,0) -- ++(0,0,-0.07) node[below,fb tick] {$\t$};
  }
  % vertical guides down to the base plane
  \draw[fb guide] (2,2,0) -- (2,2,0.95);
  \draw[fb guide] (2,0,0) -- (2,0,1);
  \draw[fb guide] (0,2,0) -- (0,2,1);
  % Glauber point and its vertical offset: below the plane, hence drawn before the translucent plane
  \node[fb gl] (G) at (2,2,1) {};
  \draw[red!60!black,line width=0.85pt,dash pattern=on 2.5pt off 1.8pt,-{Stealth[length=4pt,width=3.4pt]}]
       (2,2,1.07) -- (2,2,1.93);
  % ---------------- on-shell plane c = (a+b)/2 ----------------
  \draw[fb plane] (0,0,0) -- (2,0,1) -- (2,2,2) -- (0,2,1) -- cycle;
  \node[font=\footnotesize,text=gray!30!black,rotate=22.2] at (1.35,0.3,0.825) {在壳平面};
  % the c axis lies in front of the plane
  \draw[fb ax] (0,0,0) -- (0,0,2.8) node[above,font=\small] {$c$};
  \foreach \t in {1,2}
    \draw[gray!50!black,line width=0.5pt] (0,0,\t) -- ++(-0.045,0,0) node[left,fb tick] {$\t$};
  % ---------------- modes ----------------
  \node[fb pt,fill=black] (H)  at (0,0,0) {};
  \node[fb pt,fill=blue!55!black] (CA) at (0,2,1) {};
  \node[fb pt,fill=orange!85!black] (CB) at (2,0,1) {};
  \node[fb pt,fill=green!40!black] (US) at (2,2,2) {};
  \node[fb pt,fill=green!60!black] (S) at (1,1,1) {};
  \node[fb lab,anchor=north east] at (H.south west) {硬 $(0,0,0)$};
  \node[fb lab,anchor=east] at (CA.west) {$A$ 共线 $(0,2,1)$};
  \node[fb lab,anchor=west] at (CB.east) {$B$ 共线 $(2,0,1)$};
  \node[fb lab,anchor=west] at (US.east) {超软 $(2,2,2)$};
  \node[fb lab,anchor=south] at ($(S.north)+(0cm,0.02cm)$) {软 $(1,1,1)$};
  \node[fb lab,anchor=west,text=red!55!black] (GL) at (canvas cs:x=3.15cm,y=3.25cm) {Glauber $(2,2,1)$};
  \node[fb lab,anchor=north west,text=red!55!black,inner sep=1pt] at ($(GL.south west)+(0.1cm,0cm)$) {$k^2\approx-k_T^2$，离壳};
  \draw[red!55!black,line width=0.4pt] ($(G.east)+(0.02cm,0.03cm)$) -- (GL.west);
  % ---------------- legend ----------------
  \node[font=\small,anchor=east] at (canvas cs:x=-0.45cm,y=4.3cm) {$(k^+,\,k^-,\,k_T)\sim(\lambda^a,\,\lambda^b,\,\lambda^c)\,Q$};
  \node[font=\footnotesize,anchor=north west,align=left,inner sep=0pt] at (canvas cs:x=6.45cm,y=3.38cm)
    {平面上 $c=\tfrac12(a+b)$：\\[1pt]$\quad k_T^2\sim k^+k^-\sim|k^2|$\\[5pt]
     平面下 $c<\tfrac12(a+b)$：\\[1pt]$\quad k_T^2\gg|k^+k^-|$，$k^2\approx-k_T^2$};
\end{tikzpicture}
\caption{动量模式在指数空间 $(a,b,c)$ 中的位置，$(k^+,k^-,k_T)\sim(\lambda^a,\lambda^b,\lambda^c)\,Q$，$\lambda\ll1$（指数越大，相应分量越小）。
硬、$A$ 共线、$B$ 共线、软与超软模式都在在壳平面 $c=\frac12(a+b)$（半透明灰色平行四边形）上，即 $k_T^2$、$k^+k^-$ 与虚度 $|k^2|$ 同阶；
硬、两种共线与超软模式恰为平行四边形的四个顶点，软模式 $(1,1,1)$（TMD 因子化所用的标度）为其中心。
Glauber 模式 $(2,2,1)$（红色空心圆，透过平面看到）与超软模式有相同的 $k^\pm$，只是横动量大 $\lambda^{-1}$ 倍（红色虚线），因而位于平面之下：
$|k^+k^-|\ll k_T^2$，$k^2\approx-k_T^2$，以横动量为主且离壳；在 $(\ln k^+,\ln k^-)$ 平面（\cref{fig:momentum-regions}）上它与超软模式重合，只有加上第三个方向 $\ln k_T$ 才能区分。
本书的“超软”即 Libby--Sterman 的“软”。}
\label{fig:ir-mode-exponents}
\end{figure}
%%% END fig:ir-mode-exponents

%%% BEGIN fig:landau-pinch
\begin{figure}[htbp]
\centering
\begin{tikzpicture}[fd,
  t2 pole/.pic={\draw[line width=1pt,line cap=round] (-0.085,-0.085)--(0.085,0.085) (-0.085,0.085)--(0.085,-0.085);},
  t2 contour/.style={draw=blue!65!black,line width=1.15pt,postaction={decorate,decoration={markings,
      mark=at position #1 with {\arrow{Stealth[length=6pt,width=5pt]}}}}},
  t2 contour/.default=0.5,
  t2 move/.style={draw=gray!65!black,line width=0.6pt,dash pattern=on 2pt off 1.5pt,-{Stealth[length=4pt,width=3.4pt]}},
  t2 ax/.style={lcaxis,gray!60!black},
  t2 lab/.style={font=\footnotesize,text=black!80}]
  % ---------------- (a) poles on the same side: deformable ----------------
  \begin{scope}
    \draw[t2 ax] (-1.95,0) -- (2.0,0) node[below,xshift=-6pt] {$\Real k^-$};
    \draw[t2 ax] (0,-1.25) -- (0,1.35) node[left] {$\Imag k^-$};
    \draw[t2 contour=0.52] (-1.85,0) -- (-1.35,0) .. controls (-0.95,0) and (-0.85,0.82) .. (0,0.82)
         .. controls (0.85,0.82) and (0.95,0) .. (1.35,0) -- (1.85,0);
    \pic at (-0.55,-0.40) {t2 pole};  \pic at (0.62,-0.30) {t2 pole};
    \draw[t2 move] (-0.75,-1.05) -- (-0.59,-0.52);
    \draw[t2 move] (0.95,-1.0) -- (0.67,-0.42);
    \node[t2 lab,anchor=west] at (0.62,0.78) {形变后的围道};
    \node at (0,-1.75) {(a)\ 同侧：围道可形变避开};
  \end{scope}
  % ---------------- (b) poles on opposite sides: pinch ----------------
  \begin{scope}[xshift=5.25cm]
    \draw[t2 ax] (-1.95,0) -- (2.0,0) node[below,xshift=-6pt] {$\Real k^-$};
    \draw[t2 ax] (0,-1.25) -- (0,1.35) node[left] {$\Imag k^-$};
    \draw[t2 contour=0.8] (-1.85,0) -- (1.85,0);
    \pic at (0.16,0.2) {t2 pole};  \pic at (-0.16,-0.2) {t2 pole};
    \draw[t2 move] (0.95,0.95) -- (0.27,0.3);
    \draw[t2 move] (-0.95,-0.95) -- (-0.27,-0.3);
    \node[t2 lab,anchor=north] at (0.75,-0.18) {钳夹点};
    \node at (0,-1.75) {(b)\ 异侧相向趋近：围道被钳住};
  \end{scope}
  % ---------------- (c) endpoint ----------------
  \begin{scope}[xshift=10.5cm]
    \draw[t2 ax] (-1.95,0) -- (2.0,0) node[below,xshift=-6pt] {$\Real k^-$};
    \draw[t2 ax] (0,-1.25) -- (0,1.35) node[left] {$\Imag k^-$};
    \draw[t2 contour=0.6] (-1.1,0) .. controls (-0.95,0.55) and (-0.45,0.72) .. (0.25,0.72)
         .. controls (0.95,0.72) and (1.2,0) .. (1.85,0);
    \fill (-1.1,0) circle[radius=1.9pt];
    \node[below left=-1pt] at (-1.1,0) {$a$};
    \node[t2 lab,anchor=south east] at (-1.12,0.08) {固定端点};
    \pic at (-0.9,-0.24) {t2 pole};
    \draw[t2 move] (-0.25,-0.95) -- (-0.8,-0.33);
    \node[t2 lab,anchor=west] at (-0.2,-0.95) {趋近端点};
    \node at (0,-1.75) {(c)\ 端点：无法避开};
  \end{scope}
\end{tikzpicture}
\caption{复 $k^-$ 平面中积分围道（蓝色）的形变与 Landau 奇异性。灰色虚箭头表示外部参数趋近奇异点（或 $\ii0\to0$）时极点的移动方向。
(a)~两个极点位于实轴同侧（此处都在下半平面）：即使它们趋近实轴，围道也可以向另一侧（上半平面）形变而远离极点，积分保持解析；
(b)~两个极点从实轴两侧相向趋近同一点：围道被夹在两者之间，无论向哪一侧形变都要越过一个极点，这就是 Landau 方程所刻画的钳夹（pinch）奇异性；
(c)~极点趋近积分区间的端点 $a$：端点不能移动，围道在端点附近无法避开极点，产生端点奇异性。
只有 (b)、(c) 两种情形可能使积分作为外部参数的函数出现奇异性。多维积分中，奇异性要求积分围道（超曲面）在所有方向上都无法形变避开奇点，由此导出 Landau 方程。}
\label{fig:landau-pinch}
\end{figure}
%%% END fig:landau-pinch

%%% BEGIN fig:momentum-regions
\begin{figure}[htbp]
\centering
\begin{tikzpicture}[fd,x=2.05cm,y=2.05cm,
  t2 mode/.style={circle,draw=black,line width=0.6pt,inner sep=0pt,minimum size=6.6pt},
  t2 hyp/.style={draw=black!65,line width=0.7pt},
  t2 lab/.style={font=\footnotesize,inner sep=1.5pt},
  t2 tick/.style={font=\footnotesize},
  t2 grid/.style={draw=gray!45,line width=0.4pt,dash pattern=on 1pt off 2pt}]
  % coordinates: x = ln(k^+/Q)/ln(1/lambda), y = ln(k^-/Q)/ln(1/lambda)
  % Glauber shading for k_T ~ lambda Q:  k^+ k^- << k_T^2 ~ lambda^2 Q^2  <=>  x + y < -2
  \fill[red!7] (-3.4,-3.4) -- (0.8,-3.4) -- (0.8,-2.8) -- (-2.8,0.8) -- (-3.4,0.8) -- cycle;
  % axes, ticks, grid
  \draw[lcaxis] (-3.45,-3.45) -- (1.1,-3.45) node[below left=1pt and -4pt] {$\ln k^+$};
  \draw[lcaxis] (-3.45,-3.45) -- (-3.45,1.1) node[left] {$\ln k^-$};
  \foreach \v/\t in {-2/{\lambda^2Q},-1/{\lambda Q},0/{Q}} {
    \draw[line width=0.5pt] (\v,-3.45) -- ++(0,0.07);
    \node[t2 tick,below] at (\v,-3.45) {$\t$};
    \draw[line width=0.5pt] (-3.45,\v) -- ++(0.07,0);
    \node[t2 tick,left] at (-3.45,\v) {$\t$};
    \draw[t2 grid] (\v,-3.38) -- (\v,0.8);
    \draw[t2 grid] (-3.38,\v) -- (0.8,\v);
  }
  % constant-virtuality hyperbolas  k^+ k^- = const : straight lines of slope -1
  \draw[t2 hyp] (0.8,-0.8) -- node[pos=0.86,sloped,above=1.5pt,t2 lab] {$k^+k^-\sim Q^2$} (-0.8,0.8);
  \draw[t2 hyp] (0.8,-2.8) -- node[pos=0.36,sloped,below=2.5pt,t2 lab] {$k^+k^-\sim\lambda^2Q^2$} (-2.8,0.8);
  \draw[t2 hyp] (-0.8,-3.2) -- node[pos=0.25,sloped,above=1.5pt,t2 lab] {$k^+k^-\sim\lambda^4Q^2$} (-3.2,-0.8);
  % modes
  \node[t2 mode,fill=gray!55] (h)  at (0,0) {};
  \node[t2 mode,fill=blue!55!black] (n) at (0,-2) {};
  \node[t2 mode,fill=orange!85!black] (nb) at (-2,0) {};
  \node[t2 mode,fill=green!55!black] (s) at (-1,-1) {};
  \node[t2 mode,fill=green!25] (us) at (-2,-2) {};
  \node[t2 lab,anchor=south west] at (h.north east) {硬 $(1,1,1)$};
  \node[t2 lab,anchor=south west] at (n.north east) {$n$ 共线 $(1,\lambda^2,\lambda)$};
  \node[t2 lab,anchor=south west,fill=white] at (nb.north east) {$\bar n$ 共线 $(\lambda^2,1,\lambda)$};
  \node[t2 lab,anchor=south west,fill=white] at (s.north east) {软 $(\lambda,\lambda,\lambda)$};
  \node[t2 lab,anchor=south west,fill=red!7] at ($(us.north east)+(0.06,0.03)$) {超软 $(\lambda^2,\lambda^2,\lambda^2)$};
  \draw[red!60!black,line width=0.8pt] (-2.1,-2.1) rectangle (-1.9,-1.9);
  \node[t2 lab,anchor=north east,align=right,text=red!55!black] at ($(us.south west)+(-0.04,-0.04)$) {Glauber\\$(\lambda^2,\lambda^2,\lambda)$};
  \node[t2 lab,align=center,text=red!55!black,fill=red!7] at (0.17,-3.05) {Glauber 区\\$|k^+k^-|\ll k_T^2$};
\end{tikzpicture}
\caption{$(\ln k^+,\ln k^-)$ 平面上的动量区域（括号内为 $(k^+,k^-,k_T)/Q$ 的标度，$\lambda\ll1$；坐标轴以 $\ln(1/\lambda)$ 为单位）。$A$（$n$）共线区的大分量为 $k^+$，$B$（$\bar n$）共线区的大分量为 $k^-$。
等虚度双曲线 $k^+k^-=$ 常数在对数坐标中是斜率为 $-1$ 的直线：共线区与软区位于同一条双曲线 $k^+k^-\sim\lambda^2Q^2$ 上，彼此只差快度 $y=\frac12\ln(k^+/k^-)$；
超软区位于 $k^+k^-\sim\lambda^4Q^2$。Glauber 区由 $|k^+k^-|\ll k_T^2$ 定义：阴影为 $k_T\sim\lambda Q$ 时满足此条件的区域（通常所说的 Glauber 区是其中 $k^\pm\ll Q$ 的软部分），
其典型标度 $(\lambda^2,\lambda^2,\lambda)$ 在本图中的投影与超软区重合（红框），二者只在 $k_T$ 上不同：Glauber 动量以横动量为主，$k^2\approx-k_T^2$。}
\label{fig:momentum-regions}
\end{figure}
%%% END fig:momentum-regions

%%% BEGIN fig:reduced-dis
\begin{figure}[htbp]
\centering
\begin{tikzpicture}[fd,
  t2 hadron/.style={draw=black!70,line width=2.2pt,postaction={decorate,decoration={markings,
      mark=at position 0.55 with {\arrow[black!70]{Stealth[length=7pt,width=7.5pt]}}}}},
  t2 lab/.style={font=\small}]
  \node[hard blob,ellipse,minimum width=1.15cm,minimum height=0.78cm,inner sep=0pt] (H)  at (-2.6,0.8) {};
  \node[hard blob,ellipse,minimum width=1.15cm,minimum height=0.78cm,inner sep=0pt] (Hd) at ( 2.6,0.8) {};
  \node[jetB blob,ellipse,minimum width=3.0cm,minimum height=0.8cm] (J) at (0,2.3) {};
  \node[jet blob,ellipse,minimum width=5.0cm,minimum height=1.2cm] (A) at (0,-1.55) {};
  \node[soft blob,circle,minimum size=0.76cm,inner sep=0pt] (S) at (0,0.42) {};
  % photons
  \draw[photon] (-4.0,2.05) -- (H.135);
  \draw[photon] (Hd.45) -- (4.0,2.05);
  \path[mom={$q$},mom sep=0.28cm] (-4.0,2.05) -- (H.135);
  \path[mom={$q$},mom sep=0.28cm] (Hd.45) -- (4.0,2.05);
  % hadron lines
  \draw[t2 hadron] (-4.6,-1.55) -- (A.west);
  \draw[t2 hadron] (A.east) -- (4.6,-1.55);
  % fermion line: A -> H -> J -> H^dagger -> A
  \draw[quark] (A.165) -- (H.205);
  \draw[quark] (H.62) -- (J.195);
  \draw[quark] (J.345) -- (Hd.118);
  \draw[quark] (Hd.335) -- (A.15);
  % longitudinally polarized collinear gluons (blue): A -> H, A -> H^dagger
  \draw[longit] (A.154) -- (H.262);  \longitmark{(A.154)}{(H.262)}
  \draw[longit] (A.143) -- (H.312);  \longitmark{(A.143)}{(H.312)}
  \draw[longit] (A.30) -- (Hd.270);  \longitmark{(A.30)}{(Hd.270)}
  % soft gluons: S -> J and S -> A, on both sides of the cut
  \draw[gluon] (S.140) -- (-0.72,1.95);
  \draw[gluon] (S.40)  -- ( 0.72,1.95);
  \draw[gluon] (S.235) -- (-0.5,-0.97);
  \draw[gluon] (S.305) -- ( 0.5,-0.97);
  % cut
  \draw[cut] (0,3.0) -- (0,-2.45);
  % labels
  \node[t2 lab] at (-0.8,2.3) {$J$};
  \node[t2 lab] at (-1.15,-1.62) {$A$};
  \node[t2 lab] at (-0.62,0.42) {$S$};
  \node[t2 lab] at (H) {$H$};
  \node[t2 lab] at (Hd) {$H^\dagger$};
  \node[t2 lab,anchor=south] at (-4.35,-1.5) {$P$};
  \node[t2 lab,anchor=south] at ( 4.35,-1.5) {$P$};
  \node[font=\footnotesize,text=red!65!black,anchor=south] at (0,3.0) {末态切割};
  \node[font=\footnotesize,text=blue!55!black,align=center] at (1.2,0.4) {纵向极化\\共线胶子};
\end{tikzpicture}
\caption{深度非弹散射（DIS）的约化割图。靶强子 $P$ 的共线子图 $A$（蓝色）跨越末态切割，其中的旁观者（靶残余）穿过切割；
$H$、$H^\dagger$ 分别是振幅与共轭振幅中的硬子图，与虚光子 $q$ 相连；被击出的夸克形成末态喷注 $J$（沿 $n_B$ 方向）并穿过切割；
软子图 $S$ 以胶子与 $A$、$J$ 相连。$A$ 与 $H$（以及 $H^\dagger$）之间除一条物理极化的夸克线外，还可以有任意多条标量（纵向）极化、
即 $\epsilon^\mu\propto k^\mu$ 的共线胶子（蓝色，以圆点标记），它们在幂次计数中不受压低，经 Ward 恒等式后成为 PDF 定义中的 Wilson 线（\cref{fig:ward-identity}）。}
\label{fig:reduced-dis}
\end{figure}
%%% END fig:reduced-dis

%%% BEGIN fig:reduced-dy
\begin{figure}[htbp]
\centering
\begin{tikzpicture}[fd,
  t2 hadron/.style={draw=black!70,line width=2.2pt,postaction={decorate,decoration={markings,
      mark=at position 0.55 with {\arrow[black!70]{Stealth[length=7pt,width=7.5pt]}}}}},
  t2 lab/.style={font=\small}]
  % collinear subgraph A: amplitude part, conjugate part, spectators crossing the cut
  \node[jet blob,ellipse,minimum width=1.25cm,minimum height=0.85cm,inner sep=0pt] (Aa) at (-2.75,-1.05) {};
  \node[jet blob,ellipse,minimum width=1.25cm,minimum height=0.85cm,inner sep=0pt] (Ac) at (-2.75, 1.05) {};
  \node[jetB blob,ellipse,minimum width=1.25cm,minimum height=0.85cm,inner sep=0pt] (Ba) at ( 2.75,-1.05) {};
  \node[jetB blob,ellipse,minimum width=1.25cm,minimum height=0.85cm,inner sep=0pt] (Bc) at ( 2.75, 1.05) {};
  \node[hard blob,circle,minimum size=0.66cm,inner sep=0pt] (H)  at (0,-1.3) {};
  \node[hard blob,circle,minimum size=0.66cm,inner sep=0pt] (Hd) at (0, 1.3) {};
  \node[soft blob,circle,minimum size=0.84cm,inner sep=0pt] (S) at (0,0) {};
  % spectators of A and B crossing the cut (fermion number flows from amplitude to conjugate)
  \draw[quark=0.78] ($(Aa.north)+(-0.3,-0.04)$) -- ($(Ac.south)+(-0.3,0.04)$);
  \draw[gluon] ($(Aa.north)+(0,0.02)$) -- ($(Ac.south)+(0,-0.02)$);
  \draw[quark=0.78] ($(Aa.north)+(0.3,-0.04)$) -- ($(Ac.south)+(0.3,0.04)$);
  \draw[quark=0.78] ($(Ba.north)+(-0.3,-0.04)$) -- ($(Bc.south)+(-0.3,0.04)$);
  \draw[gluon] ($(Ba.north)+(0,0.02)$) -- ($(Bc.south)+(0,-0.02)$);
  \draw[quark=0.78] ($(Ba.north)+(0.3,-0.04)$) -- ($(Bc.south)+(0.3,0.04)$);
  % hadron lines
  \draw[t2 hadron] (-4.55,-1.65) -- (Aa.200);
  \draw[t2 hadron] (Ac.160) -- (-4.55,1.65);
  \draw[t2 hadron] (4.55,-1.65) -- (Ba.340);
  \draw[t2 hadron] (Bc.20) -- (4.55,1.65);
  % active partons:  A -> H (quark),  H -> B (antiquark: fermion arrow away from H),  conjugate side mirrored
  \draw[quark] (Aa.345) -- (H.185);
  \draw[quark] (H.355) -- (Ba.195);
  \draw[quark] (Bc.165) -- (Hd.5);
  \draw[quark] (Hd.175) -- (Ac.15);
  % virtual photon (current insertions)
  \draw[photon] (H.south) -- (0,-2.55);
  \draw[photon] (Hd.north) -- (0,2.55);
  \node[t2 lab,anchor=west] at (0.12,-2.35) {$\gamma^*(q)$};
  \node[t2 lab,anchor=west] at (0.12, 2.35) {$\gamma^*(q)$};
  % soft gluons S -> A, S -> B on both sides of the cut
  \draw[gluon] (S.200) -- (Aa.40);
  \draw[gluon] (S.160) -- (Ac.320);
  \draw[gluon] (S.340) -- (Ba.140);
  \draw[gluon] (S.20)  -- (Bc.220);
  % cut (horizontal: amplitude below, complex-conjugate amplitude above)
  \draw[cut] (-4.75,0) -- (4.75,0);
  % labels
  \node[t2 lab] at (Aa) {$A$}; \node[t2 lab] at (Ac) {$A$};
  \node[t2 lab] at (Ba) {$B$}; \node[t2 lab] at (Bc) {$B$};
  \node[t2 lab] at (H) {$H$};  \node[t2 lab] at (Hd) {$H^\dagger$};
  \node[t2 lab,anchor=south] at (0,0.44) {$S$};
  \node[t2 lab,anchor=north] at (-4.4,-1.75) {$P_A$};
  \node[t2 lab,anchor=south] at (-4.4, 1.75) {$P_A$};
  \node[t2 lab,anchor=north] at ( 4.4,-1.75) {$P_B$};
  \node[t2 lab,anchor=south] at ( 4.4, 1.75) {$P_B$};
  \node[font=\footnotesize,text=red!65!black,anchor=west] at (4.8,0) {末态切割};
  \node[font=\footnotesize,anchor=west] at (4.8,-0.75) {振幅};
  \node[font=\footnotesize,anchor=west] at (4.8, 0.75) {共轭振幅};
\end{tikzpicture}
\caption{Drell--Yan 过程 $A+B\to\gamma^*+X$ 的约化割图（强子张量）。为使两个强子分列左右，切割（红色虚线）画成水平的：下半部为振幅，上半部为共轭振幅。
强子 $A$、$B$ 的共线子图（蓝、橙色）各分为振幅与共轭振幅两部分，二者之间的旁观者（强子残余）线穿过切割；硬子图 $H$、$H^\dagger$ 各由一条来自 $A$ 的夸克线与一条来自 $B$ 的反夸克线相连，
并通过电流插入与 $\gamma^*$ 相连；软子图 $S$ 以胶子与 $A$、$B$ 相连，在切割两侧均可附着。共线子图与 $H$ 之间的纵向极化胶子未画出（参见\cref{fig:reduced-dis}）；
旁观者之间的 Glauber 交换见\cref{fig:glauber-exchange}。}
\label{fig:reduced-dy}
\end{figure}
%%% END fig:reduced-dy

%%% BEGIN fig:reduced-ee
\begin{figure}[htbp]
\centering
\begin{tikzpicture}[fd,
  t2 lab/.style={font=\small}]
  \node[hard blob,circle,minimum size=0.6cm,inner sep=0pt] (H)  at (-2.75,0) {};
  \node[hard blob,circle,minimum size=0.6cm,inner sep=0pt] (Hd) at ( 2.75,0) {};
  \node[jet blob,ellipse,minimum width=3.3cm,minimum height=0.86cm] (J1) at (0, 1.6) {};
  \node[jetB blob,ellipse,minimum width=3.3cm,minimum height=0.86cm] (J2) at (0,-1.6) {};
  \node[soft blob,circle,minimum size=0.8cm,inner sep=0pt] (S) at (0,0) {};
  % photons
  \draw[photon] (-4.25,0) -- (H.west);
  \draw[photon] (Hd.east) -- (4.25,0);
  \path[mom={$q$},mom sep=0.28cm] (-4.25,0) -- (H.west);
  \path[mom={$q$},mom sep=0.28cm] (Hd.east) -- (4.25,0);
  % fermion lines: H -> J1 -> H^dagger -> J2 -> H
  \draw[quark] (H.65) -- (J1.188);
  \draw[quark] (J1.352) -- (Hd.115);
  \draw[quark] (Hd.245) -- (J2.8);
  \draw[quark] (J2.172) -- (H.295);
  % soft gluons (to both jets, on both sides of the cut)
  \draw[gluon] (S.152) -- (-0.92,1.2);
  \draw[gluon] (S.28) -- (0.92,1.2);
  \draw[gluon] (S.208) -- (-0.92,-1.2);
  \draw[gluon] (S.332) -- (0.92,-1.2);
  % cut
  \draw[cut] (0,2.45) -- (0,-2.45);
  % labels
  \node[t2 lab] at (-0.85,1.6) {$J_1$};
  \node[t2 lab] at (-0.85,-1.6) {$J_2$};
  \node[t2 lab,anchor=east] at ($(S.west)+(-0.04,0)$) {$S$};
  \node[t2 lab,anchor=north east] at (H.south west) {$H$};
  \node[t2 lab,anchor=north west] at (Hd.south east) {$H^\dagger$};
  \node[font=\footnotesize,anchor=west] at (1.75,1.95) {沿 $p_1$};
  \node[font=\footnotesize,anchor=west] at (1.75,-1.95) {沿 $p_2$};
  \node[font=\footnotesize,text=red!65!black,anchor=south] at (0,2.45) {末态切割};
\end{tikzpicture}
\caption{$e^+e^-\to$ 两喷注的约化割图（领头钳夹面）。$H$（振幅）与 $H^\dagger$（共轭振幅）是全部离壳线收缩成的硬点；$J_1$、$J_2$
是分别沿类光方向 $p_1$、$p_2$ 运动的共线子图（喷注），各以一条物理极化的夸克线与 $H$、$H^\dagger$ 相连，并穿过末态切割（红色虚线）；
软子图 $S$ 只通过胶子与两个喷注相连，不与硬部分相连，其末态粒子同样穿过切割。在 Feynman 规范下喷注与硬部分之间还可以有任意多条标量（纵向）极化的共线胶子，
图中未画出（参见\cref{fig:reduced-dis}）。}
\label{fig:reduced-ee}
\end{figure}
%%% END fig:reduced-ee

%%% BEGIN fig:vertex-regions
\begin{figure}[htbp]
\centering
\begin{tikzpicture}[fd,
  t2 soft/.style={draw=green!45!black},
  t2 jetA/.style={draw=blue!55!black,fill=blue!10,line width=0.5pt},
  t2 jetB/.style={draw=orange!75!black,fill=orange!14,line width=0.5pt},
  t2 hard/.style={draw=black!70,fill=gray!35,line width=0.6pt},
  t2 lab/.style={font=\footnotesize}]
  % common geometry:  V = photon vertex,  E1/E2 = ends of the q / qbar legs,  G1/G2 = gluon vertices
  % ======================= (a) soft gluon =======================
  \begin{scope}
    \coordinate (V) at (0,0); \coordinate (E1) at (28:2.75); \coordinate (E2) at (-28:2.75);
    \coordinate (G1) at ($(V)!0.62!(E1)$); \coordinate (G2) at ($(V)!0.62!(E2)$);
    \begin{scope}[on background layer]
      \draw[t2 jetA,rotate around={28:($(V)!0.52!(E1)$)}] ($(V)!0.52!(E1)$) ellipse[x radius=1.55cm,y radius=0.27cm];
      \draw[t2 jetB,rotate around={-28:($(V)!0.52!(E2)$)}] ($(V)!0.52!(E2)$) ellipse[x radius=1.55cm,y radius=0.27cm];
    \end{scope}
    \draw[photon] (-1.15,0) -- (V);
    \draw[quark] (V) -- (G1);  \draw[quark] (G1) -- (E1);
    \draw[quark] (E2) -- (G2); \draw[quark] (G2) -- (V);
    \draw[gluon,t2 soft] (G2) -- (G1);
    \filldraw[t2 hard] (V) circle[radius=0.15];
    \foreach \p in {G1,G2} \node[vertex] at (\p) {};
    \path[mom={$q$},mom sep=0.27cm] (-1.15,0) -- (-0.15,0);
    \path[mom={$p_1$},mom pos=0.6] (G1) -- (E1);
    \path[mom'={$p_2$},mom pos=0.6] (G2) -- (E2);
    \path[mom'={$k$},mom sep=0.3cm] (G2) -- (G1);
    \node[t2 lab,text=green!40!black,anchor=east] at ($(G1)!0.5!(G2)+(-0.16,0)$) {$S$};
    \node at (0.75,-2.05) {(a)\ 软区：$k\to0$};
  \end{scope}
  % ======================= (b) gluon collinear to q =======================
  \begin{scope}[xshift=5.2cm]
    \coordinate (V) at (0,0); \coordinate (E1) at (28:2.75); \coordinate (E2) at (-28:2.75);
    \coordinate (G1) at ($(V)!0.8!(E1)$); \coordinate (G2) at ($(V)!0.12!(E2)$);
    \begin{scope}[on background layer]
      \draw[t2 jetA,rotate around={31:($(V)!0.5!(E1)+(0.05,-0.12)$)}] ($(V)!0.5!(E1)+(0.05,-0.12)$) ellipse[x radius=1.6cm,y radius=0.5cm];
      \draw[t2 jetB,rotate around={-28:($(V)!0.62!(E2)$)}] ($(V)!0.62!(E2)$) ellipse[x radius=1.15cm,y radius=0.24cm];
    \end{scope}
    \draw[quark=0.35] (V) -- (G1);  \draw[quark] (G1) -- (E1);
    \draw[quark] (E2) -- (G2); \draw[plain] (G2) -- (V);
    \draw[gluon] (G2) .. controls ($(V)!0.42!(E1)+(0.22,-0.5)$) and ($(V)!0.62!(E1)+(0.12,-0.38)$) .. (G1);
    \filldraw[t2 hard] ($(V)!0.5!(G2)$) circle[radius=0.33];
    \draw[photon] (-1.15,0) -- ($(V)!0.5!(G2)+(-0.33,0.03)$);
    \foreach \p in {G1} \node[vertex] at (\p) {};
    \node[t2 lab] at ($(V)!0.5!(G2)$) {$H$};
    \path[mom={$q$},mom sep=0.27cm] (-1.15,0) -- (-0.3,0);
    \path[mom={$p_1$},mom pos=0.55] (G1) -- (E1);
    \path[mom'={$p_2$},mom pos=0.6] (G2) -- (E2);
    \path[mom'={$k$},mom sep=0.3cm] (G2) .. controls ($(V)!0.42!(E1)+(0.22,-0.5)$) and ($(V)!0.62!(E1)+(0.12,-0.38)$) .. (G1);
    \node[t2 lab,text=blue!50!black] at (0.95,1.25) {$J_1$};
    \node at (0.75,-2.05) {(b)\ $k\parallel p_1$};
  \end{scope}
  % ======================= (c) gluon collinear to qbar =======================
  \begin{scope}[xshift=10.4cm]
    \coordinate (V) at (0,0); \coordinate (E1) at (28:2.75); \coordinate (E2) at (-28:2.75);
    \coordinate (G1) at ($(V)!0.12!(E1)$); \coordinate (G2) at ($(V)!0.8!(E2)$);
    \begin{scope}[on background layer]
      \draw[t2 jetB,rotate around={-31:($(V)!0.5!(E2)+(0.05,0.12)$)}] ($(V)!0.5!(E2)+(0.05,0.12)$) ellipse[x radius=1.6cm,y radius=0.5cm];
      \draw[t2 jetA,rotate around={28:($(V)!0.62!(E1)$)}] ($(V)!0.62!(E1)$) ellipse[x radius=1.15cm,y radius=0.24cm];
    \end{scope}
    \draw[plain] (V) -- (G1);  \draw[quark] (G1) -- (E1);
    \draw[quark] (E2) -- (G2); \draw[quark=0.65] (G2) -- (V);
    \draw[gluon] (G1) .. controls ($(V)!0.42!(E2)+(0.22,0.5)$) and ($(V)!0.62!(E2)+(0.12,0.38)$) .. (G2);
    \filldraw[t2 hard] ($(V)!0.5!(G1)$) circle[radius=0.33];
    \draw[photon] (-1.15,0) -- ($(V)!0.5!(G1)+(-0.33,-0.03)$);
    \foreach \p in {G2} \node[vertex] at (\p) {};
    \node[t2 lab] at ($(V)!0.5!(G1)$) {$H$};
    \path[mom={$q$},mom sep=0.27cm] (-1.15,0) -- (-0.3,0);
    \path[mom={$p_1$},mom pos=0.6] (G1) -- (E1);
    \path[mom'={$p_2$},mom pos=0.55] (G2) -- (E2);
    \path[mom={$-k$},mom sep=0.3cm] (G1) .. controls ($(V)!0.42!(E2)+(0.22,0.5)$) and ($(V)!0.62!(E2)+(0.12,0.38)$) .. (G2);
    \node[t2 lab,text=orange!60!black] at (0.95,-1.25) {$J_2$};
    \node at (0.75,-2.05) {(c)\ $k\parallel p_2$};
  \end{scope}
\end{tikzpicture}
\caption{单圈无质量顶点 $\gamma^*(q)\to q(p_1)\,\bar q(p_2)$ 的三个钳夹面（领头区域）。胶子动量 $k$ 由反夸克线流向夸克线，
内夸克线携带 $p_1-k$，内反夸克线携带 $p_2+k$。
(a)~软区 $k\to0$：两条内费米子线在壳并分别平行于 $p_1$、$p_2$（蓝、橙色区域为两个喷注），零动量胶子（绿色）构成软子图 $S$，硬部分只剩光子顶点；
(b)~$k=\alpha p_1$（$0<\alpha<1$）：胶子与内夸克线都在壳且平行于 $p_1$，同属喷注 $J_1$，而内反夸克线的虚度 $(p_2+\alpha p_1)^2=\alpha q^2$ 不小，收缩进硬部分 $H$；
(c)~$k=-\beta p_2$（$0<\beta<1$）：与 (b) 镜像对称，胶子携带物理动量 $-k=\beta p_2$ 沿 $p_2$ 方向运动并属于 $J_2$，内夸克线的虚度 $(p_1+\beta p_2)^2=\beta q^2$ 不小。
费米子线上的箭头表示费米子数流向，动量箭头单独画出。}
\label{fig:vertex-regions}
\end{figure}
%%% END fig:vertex-regions
